% ============================================================
%   CHAPTER 4: Description of the Design Elements
% ============================================================
\chapter{Description of the Design Elements}
\label{chap:design_elements}

\begin{table}[H]
\caption{Revision History – Chapter 4}
\centering
\begin{tabular}{llll}
\toprule
Version & Date & Author & Description \\
\midrule
1.0 & \today & Y. Sangale & Initial Draft \\
\bottomrule
\end{tabular}
\end{table}

\textcolor{blue}{\textbf{\textit{This Section will be completed later.}}}

This chapter provides a detailed technical description of all major design elements composing the \gls{rv64i}
64-bit single-core microcontroller.  
Each sub-block is presented in terms of hardware interfaces, bus connectivity, internal registers, 
and its functional and structural roles within the \gls{SoC}.  
All diagrams and tables are placeholders for later detailed schematics.

% ============================================================
\section{\gls{cpu} Core}

\subsection{HW Interfaces to Other Blocks}
The \gls{cpu} Core is the central processing unit of the system, implementing the RV64IM instruction set
and acting as the single bus master in the microcontroller architecture.  
It connects to all major subsystems as follows:
\begin{itemize}
    \item \textbf{Instruction Memory (\gls{I-MEM}):} Connected directly via the \textbf{Instruction Bus (I-Bus)} 
          for program fetch under the \gls{Harvard} architecture model.
    \item \textbf{Data Memory (\gls{D-MEM}):} Connected directly via the \textbf{Data Bus (D-Bus)} for load/store access.
    \item \textbf{Wishbone Bus:} Acts as the \textbf{bus master} interfacing to all memory-mapped peripherals 
          (\gls{gpio}, \gls{uart}, \gls{qspi}, Interrupt Controller, and Clock Unit).
    \item \textbf{Interrupt Input:} Receives a single aggregated interrupt signal from the Vectored Interrupt Controller (VIC).
    \item \textbf{Clock and Reset:} Driven by the system clock from the Clock Generation Unit (CGU) and 
          synchronized reset signal.
\end{itemize}

\begin{figure}[H]
    \centering
    \caption{\gls{cpu} Core interface diagram showing I-Bus, D-Bus, and Wishbone connections to system blocks.}
    \label{fig:cpu_interface}
\end{figure}

\subsection{Bus Interface}
The \gls{cpu} operates as a \textbf{Wishbone Master} for all peripheral transactions.
It initiates read/write cycles with standard Wishbone control signals:
\texttt{ADR\_O}, \texttt{DAT\_O}, \texttt{DAT\_I}, \texttt{STB\_O}, \texttt{ACK\_I}, \texttt{CYC\_O}, and \texttt{WE\_O}.  
Instruction and data memories are accessed directly via the independent I-Bus and D-Bus, each 
64-bit wide for full datapath utilization.

\begin{figure}[H]
    \centering
    \caption{\gls{cpu} bus architecture — direct \gls{Harvard}-style connections to \gls{I-MEM}/\gls{D-MEM} and Wishbone master for peripherals.}
    \label{fig:cpu_bus_arch}
\end{figure}

\subsection{Registers}
The \gls{cpu} core implements 32 × 64-bit general-purpose registers (\texttt{x0–x31}) and a 64-bit program counter.
Several internal control and status registers are included for interrupt and \gls{pipeline} management.

\begin{table}[H]
    \centering
    \caption{\gls{cpu} Core Register Summary}
    \label{tab:cpu_registers}
    \begin{tabular}{|l|l|}
        \hline
        \textbf{Register} & \textbf{Description} \\ \hline
        \gls{pc} & 64-bit Program Counter for sequential and branch control \\ \hline
        IR & Instruction Register holding current fetched instruction \\ \hline
        SR & Status Register with control and flag bits \\ \hline
        IBASE & Base address for interrupt vector table \\ \hline
        MSTATUS & Machine Status Register (\gls{rv64i} subset) \\ \hline
    \end{tabular}
\end{table}

\subsection{Functional Overview}
The \gls{cpu} executes instructions using a deterministic scalar \gls{pipeline} that supports the \gls{riscv} RV64IM base instruction set.
Internally, it consists of the following major sub-modules:
\begin{itemize}
    \item \textbf{Instruction Decoder:} Decodes opcodes fetched from \gls{I-MEM} and generates control signals for \gls{alu} and memory.
    \item \textbf{Program Counter (\gls{pc} ):} Maintains the address of the current instruction and computes sequential or branch targets.
    \item \textbf{\gls{alu}:} Performs arithmetic, logical, shift, and comparison operations; supports integer multiply/divide.
    \item \textbf{Register File:} Provides two read ports and one write port for operand access and result storage.
    \item \textbf{Interrupt Input:} Interfaces with VIC to handle asynchronous interrupt requests.
\end{itemize}

Each instruction is executed in a simple in-order flow without speculation or \gls{pipeline} hazards beyond single-cycle stalls.

\begin{figure}[H]
    \centering
    \caption{\gls{cpu} Core internal structure showing Instruction Decoder, \gls{alu}, Program Counter, and Register File.}
    \label{fig:cpu_structure}
\end{figure}

\subsection{Structural Overview}
The \gls{cpu} architecture follows a \gls{Harvard} model with:
\begin{itemize}
    \item Dedicated instruction path (I-Bus) for fetching 64-bit instructions.
    \item Dedicated data path (D-Bus) for memory and peripheral data transactions.
    \item Unified Wishbone interface for peripheral control and configuration.
\end{itemize}

The internal control path coordinates fetch, decode, execute, and memory operations.
Instruction flow synchronization ensures deterministic timing suitable for embedded applications.

\begin{figure}[H]
    \centering
    \caption{Simplified \gls{Harvard} architecture connection of \gls{cpu} to \gls{I-MEM}, \gls{D-MEM}, and Wishbone peripherals.}
    \label{fig:cpu_harvard_arch}
\end{figure}

\subsection{Interrupts}
A single interrupt input line connects to the VIC.  
The \gls{cpu} jumps to a vector address computed as:
\[
\text{ISR Vector} = \text{IBASE} + (\text{\gls{IRQ}\_ID} \times 0x10)
\]
Upon completion, control returns to the instruction following the interrupted instruction.

\subsection{Clock and Reset}
The \gls{cpu} operates at an 80\,MHz clock supplied by the CGU.  
Reset initializes the \gls{pc} to 0x0000\_0000, clears all \gls{pipeline} stages, and resets the internal status registers.

\subsection{Voltage Class}
Core operates in the 3.3\,V logic domain.

